%%═══════════════════════════════════════════════════════════════════
%% Lecture 04 — Perturbation Theory for the Harmonic Oscillator:
%%              Transition Amplitudes, Fermi's Golden Rule,
%%              Second-Order Perturbation Theory, and
%%              Introduction to Feynman Diagrams
%% 77609 — Introduction to Elementary Particles
%% Date: 19/04/2026
%% Lecturer: Prof. Yonit Hochberg
%%═══════════════════════════════════════════════════════════════════

\section{Lecture 04 — Perturbation Theory and Feynman Diagrams for the Harmonic Oscillator}\label{sec:lec04:perturbation}
\begin{center}
\begin{tabular}{ll}
  \textbf{Date:}\quad 19/04/2026   & \\
\end{tabular}
\end{center}
\hrule\vspace{1.5em}
%%──────────────────────────────────────────────────────────────────
\subsection*{Overview}

This lecture introduces \textbf{perturbation theory} in the context of coupled harmonic oscillators — a tractable system where results can be verified explicitly.  The central objects are:
\begin{enumerate}
  \item \textbf{Fermi's Golden Rule}: the transition rate between states expressed in terms of a transition \emph{amplitude} and a \emph{phase-space} factor.
  \item \textbf{Selection rules}: which transitions are kinematically and dynamically allowed for a given interaction Hamiltonian.
  \item \textbf{Second-order perturbation theory}: explicit computation of a $y \to 3x$ decay amplitude via a sum over intermediate states.
  \item \textbf{Feynman diagrams}: a pictorial — and ultimately computational — shorthand that encodes the same information with far less effort.
\end{enumerate}

The motivating philosophy is to \emph{prove} that Feynman diagrams work on a familiar system (the harmonic oscillator), so that one can later \emph{trust} the same machinery when applied to quantum field theory.

%%──────────────────────────────────────────────────────────────────
\subsection*{Setup: Coupled Harmonic Oscillators}

Consider a two-dimensional system with potential
\begin{equation}
  V(x,y)
  = \frac{k\,x^2}{2}
    + \frac{4k^2 y^2}{2}
    + \alpha\, x^2 y.
  \label{eq:lec04:potential}
\end{equation}

\begin{itemize}
  \item The first two terms describe two \emph{decoupled} harmonic oscillators: an $x$-mode with angular frequency $\omega_x = \sqrt{k}$ and a $y$-mode with $\omega_y = 2\omega_x$.
  \item The third term $\alpha x^2 y$ is an \emph{anharmonic} (interaction) term that couples the two modes and transfers energy between them.
  \item The coupling constant $\alpha$ is assumed small, so this term is treated as a perturbation.
\end{itemize}

\begin{remark}[Physical picture of the interaction term]
The anharmonic term $\alpha x^2 y$ can be thought of as an \emph{interaction vertex}: it can annihilate one $y$-quantum and create two $x$-quanta (or vice versa).  This is the harmonic-oscillator analogue of a particle decay in QFT.  The parameter $\alpha$ controls the \emph{rate} of energy transfer between the modes — just as a coupling constant controls an interaction rate in field theory.
\end{remark}

%%──────────────────────────────────────────────────────────────────
\subsection*{Fermi's Golden Rule}

\dfn{Fermi's Golden Rule}{The transition \emph{rate} $W$ from an initial state $|i\rangle$ to a final state $|f\rangle$ is
\begin{equation}
  W_{i \to f}
  \;\propto\;
  |\mathcal{A}|^2 \times \text{(phase space)},
  \label{eq:lec04:FGR}
\end{equation}
where $\mathcal{A} = \langle f | H_1 | i \rangle$ is the \textbf{transition amplitude} (sometimes written $\mathcal{M}$) and the phase-space factor encodes the kinematic constraints.}

\begin{itemize}
  \item The amplitude $\mathcal{A}$ contains \textbf{all dynamical information} about the interaction.  It is computed via perturbation theory or, equivalently, via Feynman diagrams.
  \item The \textbf{phase-space} factor accounts for kinematics: how much ``room'' is available in the space of final states.  For example, if a heavy particle decays into products whose combined mass is close to its own mass, the phase space is small and the decay rate is suppressed, regardless of the coupling strength.
\end{itemize}

%%──────────────────────────────────────────────────────────────────
\subsection*{Ladder Operators and Matrix Elements}

In the quantum harmonic oscillator with coordinate $q$ (standing for $x$ or $y$ or $z$), the position operator is expressed in terms of creation and annihilation operators as
\begin{equation}
  q
  = \frac{1}{\sqrt{2\omega_q}}
    \bigl(a_q + a_q^\dagger\bigr),
  \label{eq:lec04:q-ladder}
\end{equation}
where $\omega_q$ is the frequency of that mode.  The standard ladder-operator actions are
\begin{equation}
  a_q^\dagger |n\rangle = \sqrt{n+1}\,|n+1\rangle,
  \qquad
  a_q |n\rangle = \sqrt{n}\,|n-1\rangle.
  \label{eq:lec04:ladder-actions}
\end{equation}

%%──────────────────────────────────────────────────────────────────
\subsection*{Selection Rules from the Interaction Hamiltonian}

The interaction Hamiltonian for the two-oscillator system is
\begin{equation}
  H_1 = \alpha\, x^2 y.
  \label{eq:lec04:H1-two}
\end{equation}
Using \eqref{eq:lec04:q-ladder}, $x^2 \propto (a_x + a_x^\dagger)^2$ and $y \propto (a_y + a_y^\dagger)$.  Expanding these products reveals which changes in occupation numbers $\Delta n_x$ and $\Delta n_y$ are possible:

\dfn{Selection rules for $H_1 = \alpha x^2 y$}{A transition $|i\rangle \to |f\rangle$ has a non-zero amplitude \emph{only} if
\begin{equation}
  \Delta n_y = \pm 1
  \qquad \text{and} \qquad
  \Delta n_x \in \{0,\, \pm 2\}.
  \label{eq:lec04:selection-rules}
\end{equation}}

\begin{remark}[Generalisation]
For a general interaction $H_1 \propto x^p y^q$, the selection rules become $\Delta n_x \in \{-p, -p+2, \ldots, p\}$ and $\Delta n_y \in \{-q, -q+2, \ldots, q\}$.  For example, $H_1 \propto x^2 y^3$ gives $\Delta n_x \in \{0, \pm 2\}$ and $\Delta n_y \in \{\pm 1, \pm 3\}$.
\end{remark}

%%──────────────────────────────────────────────────────────────────
\subsection*{First Example: $y \to 2x$ Decay (Two-Oscillator System)}

\subsubsection*{Kinematic setup}
With $\omega_y = 2\omega_x$, energy conservation requires that the final state has $n_y = 0$ and $n_x = 2$.  Starting from $|i\rangle = |n_x, n_y\rangle = |0,1\rangle$, the allowed final state is $|f\rangle = |2,0\rangle$.

\subsubsection*{First-order amplitude}
\begin{align}
  \mathcal{A}
  &= \langle 2,0 | \,\alpha\, x^2 y\, | 0,1 \rangle \notag\\
  &= \alpha\,\langle 2,0 |
      \frac{(a_x + a_x^\dagger)^2}{2\omega_x}
      \cdot
      \frac{(a_y + a_y^\dagger)}{\sqrt{2\omega_y}}
    | 0,1 \rangle.
  \label{eq:lec04:amp-2osc}
\end{align}
The $y$ factor $a_y + a_y^\dagger$ acting on $|n_y = 1\rangle$ gives $a_y|1\rangle = |0\rangle$ (the relevant term, since $a_y^\dagger|1\rangle = \sqrt{2}|2\rangle$ does not lead to $|0\rangle$).  The $x$ factor $(a_x + a_x^\dagger)^2$ acting on $|0\rangle$ produces $|2\rangle$ via two applications of $a_x^\dagger$:
\begin{equation}
  (a_x^\dagger)^2 |0\rangle = \sqrt{2}\,|2\rangle.
  \label{eq:lec04:x-create-two}
\end{equation}
This process is described pictorially as:

\begin{center}
\emph{One $y$-particle enters the vertex and is annihilated; two $x$-particles are created and exit.}
\end{center}

\subsubsection*{The Feynman diagram}
\begin{center}
  \begin{tikzpicture}
    \begin{feynman}
      \vertex (v);
      \vertex [above left=1.5cm of v] (yi) {\(y\)};
      \vertex [above right=1.2cm of v] (xf1) {\(x\)};
      \vertex [right=1.8cm of v] (xf2) {\(x\)};
      \diagram*{
        (yi) -- [scalar, edge label=\(y\)] (v),
        (v)  -- [scalar, edge label=\(x\)] (xf1),
        (v)  -- [scalar, edge label=\(x\)] (xf2),
      };
    \end{feynman}
  \end{tikzpicture}
\end{center}

The vertex carries a factor $\alpha$ (the coupling constant).

\subsubsection*{Decay width and lifetime}
\begin{equation}
  \DecayRate \equiv \Gamma_{y \to 2x}
  \;\propto\; |\mathcal{A}|^2 \times \text{phase space}
  \;\propto\; \alpha^2.
  \label{eq:lec04:decay-width}
\end{equation}
The \textbf{lifetime} of the $y$ particle is
\begin{equation}
  \tau = \frac{1}{\Gamma} \;\propto\; \frac{1}{\alpha^2}.
  \label{eq:lec04:lifetime}
\end{equation}

%%──────────────────────────────────────────────────────────────────
\subsection*{Second Example: $y \to 3x$ Decay — Three-Oscillator System}

\subsubsection*{System and Hamiltonian}
Now consider a \emph{three-mode} system with coordinates $x$, $y$, $z$ and frequencies $\omega_x$, $\omega_y = 3\omega_x$, $\omega_z \gg \omega_x, \omega_y$.  The interaction Hamiltonian has two terms:
\begin{equation}
  H_1 = \lambda_1\, xyz + \lambda_2\, x^2 z,
  \label{eq:lec04:H1-three}
\end{equation}
where $\lambda_1$ and $\lambda_2$ are small coupling constants treated perturbatively.

States are labelled $|n_x, n_y, n_z\rangle$.  We wish to compute the transition amplitude for
\begin{equation}
  |i\rangle = |0,1,0\rangle
  \;\longrightarrow\;
  |f\rangle = |3,0,0\rangle.
  \label{eq:lec04:process}
\end{equation}
Energy conservation: $E_i = \omega_y = 3\omega_x = E_f$. \checkmark

\subsubsection*{Why second-order perturbation theory is needed}
Neither $\lambda_1 xyz$ nor $\lambda_2 x^2 z$ alone can connect $|0,1,0\rangle$ to $|3,0,0\rangle$ in a single step.  Each interaction changes occupation numbers by limited amounts (e.g.\ $\lambda_1$ changes $n_x, n_y, n_z$ by $\pm 1$ each; $\lambda_2$ changes $n_x$ by $\pm 2$ and $n_z$ by $\pm 1$).  Therefore, one must act with $H_1$ \emph{twice} (second order), passing through an intermediate state $|n\rangle$.

\dfn{Second-order perturbation theory amplitude}{\begin{equation}
  \mathcal{A}^{(2)}_{i \to f}
  = \sum_{n} \frac{\langle f | H_1 | n \rangle \langle n | H_1 | i \rangle}{E_i - E_n},
  \label{eq:lec04:2nd-order}
\end{equation}
where the sum runs over all intermediate states $|n\rangle$ with $E_n \neq E_i$.}

\begin{remark}[Virtual intermediate states]
The intermediate state $|n\rangle$ in \eqref{eq:lec04:2nd-order} need \emph{not} satisfy energy conservation ($E_n \neq E_i$ is generically the case).  These are called \textbf{virtual} (or \emph{off-shell}) intermediate states.  Only the initial and final states must be on-shell (energy-conserving).  This is the quantum mechanical precursor to virtual particles in QFT.
\end{remark}

\subsubsection*{Identifying intermediate states}

We must find all $|n\rangle$ such that both $\langle n | H_1 | i \rangle \neq 0$ and $\langle f | H_1 | n \rangle \neq 0$.

\paragraph{Intermediate state 1: $|n_1\rangle = |1,0,1\rangle$}
\begin{itemize}
  \item Step 1 (act on $|i\rangle = |0,1,0\rangle$ with $\lambda_1 xyz$): annihilate the $y$ and create one $x$ and one $z$, reaching $|1,0,1\rangle$.
  \item Step 2 (act on $|1,0,1\rangle$ with $\lambda_2 x^2 z$): annihilate the $z$ and create two more $x$, reaching $|3,0,0\rangle = |f\rangle$.
\end{itemize}

\paragraph{Intermediate state 2: $|n_2\rangle = |2,1,1\rangle$}
\begin{itemize}
  \item Step 1 (act on $|0,1,0\rangle$ with $\lambda_2 x^2 z$): create two $x$ and one $z$; $y$ is untouched, reaching $|2,1,1\rangle$.
  \item Step 2 (act on $|2,1,1\rangle$ with $\lambda_1 xyz$): annihilate the $y$ and the $z$, create one more $x$, reaching $|3,0,0\rangle = |f\rangle$.
\end{itemize}

These are the \emph{only} two contributing intermediate states.

\subsubsection*{Matrix elements}

Using $q = \frac{1}{\sqrt{2\omega_q}}(a_q + a_q^\dagger)$ and defining the shorthand
\begin{equation}
  C_1 \equiv \frac{1}{(2)^{3/2}\sqrt{\omega_x\,\omega_y\,\omega_z}},
  \qquad
  C_2 \equiv \frac{1}{(2)^{3/2}\,\omega_x\sqrt{\omega_z}},
  \label{eq:lec04:C1-C2}
\end{equation}
the relevant matrix elements evaluate to:

\begin{align}
  \langle 1,0,1 | \lambda_1 xyz | 0,1,0 \rangle
  &= \lambda_1\, C_1,
  \label{eq:lec04:me1a}\\
  \langle 3,0,0 | \lambda_2 x^2 z | 1,0,1 \rangle
  &= \lambda_2\, C_2 \cdot \sqrt{6},
  \label{eq:lec04:me1b}\\
  \langle 2,1,1 | \lambda_2 x^2 z | 0,1,0 \rangle
  &= \lambda_2\, C_2 \cdot \sqrt{2},
  \label{eq:lec04:me2a}\\
  \langle 3,0,0 | \lambda_1 xyz | 2,1,1 \rangle
  &= \lambda_1\, C_1 \cdot \sqrt{6}.
  \label{eq:lec04:me2b}
\end{align}

The $\sqrt{6}$ and $\sqrt{2}$ factors arise from the ladder-operator actions $a^\dagger|n\rangle = \sqrt{n+1}|n+1\rangle$.

\subsubsection*{Full second-order amplitude}

The energy denominators for the two channels are:
\begin{align}
  E_i - E_{n_1}
  &= 3\omega_x - (\omega_x + \omega_z)
   = 2\omega_x - \omega_z,\\
  E_i - E_{n_2}
  &= 3\omega_x - (2\omega_x + \omega_y + \omega_z)
   = -(2\omega_x + \omega_z)
   \quad [\text{using }\omega_y = 3\omega_x].
\label{eq:lec04:energy-denominators}
\end{align}

Substituting into \eqref{eq:lec04:2nd-order}:
\begin{align}
  \mathcal{A}^{(2)}
  &= \frac{(\lambda_2 C_2 \sqrt{6})(\lambda_1 C_1)}{2\omega_x - \omega_z}
   + \frac{(\lambda_1 C_1 \sqrt{6})(\lambda_2 C_2\sqrt{2})}
          {-(2\omega_x + \omega_z)}.
  \label{eq:lec04:amp-full}
\end{align}

Factoring out $\lambda_1 \lambda_2 C_1 C_2 \sqrt{6}$:
\begin{align}
  \mathcal{A}^{(2)}
  &= \lambda_1 \lambda_2\, C_1 C_2\,\sqrt{6}
    \left[
      \frac{1}{2\omega_x - \omega_z}
      - \frac{\sqrt{2}}{\omega_z + 2\omega_x}
    \right].
  \label{eq:lec04:amp-factored}
\end{align}

Taking a common denominator $(\omega_z^2 - 4\omega_x^2)$:
\begin{equation}
  \boxed{
    \mathcal{A}^{(2)}
    = -\frac{\lambda_1\lambda_2\sqrt{2}}{4\,\omega_x^2}
      \cdot\frac{1}{\omega_z^2/(4\omega_x^2) - 1}
  }
  \label{eq:lec04:amp-final}
\end{equation}
(after substituting $\omega_y = 3\omega_x$ and simplifying).  This is the complete first-nontrivial-order transition amplitude for $y \to 3x$.

%%──────────────────────────────────────────────────────────────────
\subsection*{Introduction to Feynman Diagrams}

\subsubsection*{Motivation}

The second-order calculation above requires:
\begin{enumerate}
  \item Identifying all contributing intermediate states by inspection.
  \item Computing each matrix element using ladder operators.
  \item Combining all terms with the correct energy denominators.
\end{enumerate}
For homework problems with many interactions or many modes, this becomes \emph{very} cumbersome.  Feynman diagrams provide a systematic pictorial method that automatically generates the correct answer.

\subsubsection*{Elements of a Feynman diagram}

\dfn{External states}{\textbf{External states} (also called \emph{external lines}) represent the physical, on-shell initial and final particles.  In the $y \to 3x$ example there is one external initial line ($y$) and three external final lines ($x$, $x$, $x$).}

\dfn{Vertices}{\textbf{Vertices} are points where lines meet; they correspond to interaction terms in $H_1$.  Each vertex contributes a factor of the coupling constant ($\lambda_1$ or $\lambda_2$).}

\dfn{Internal lines / propagators}{\textbf{Internal lines} (also called \emph{propagators} or \emph{virtual lines}) connect two vertices and represent the virtual intermediate state.  They correspond to the energy denominator $1/(E_i - E_n)$ in second-order perturbation theory.  Unlike external lines, internal lines need not be on-shell.}

\subsubsection*{Feynman diagram for $y \to 3x$}

The process requires one $\lambda_1$ vertex (the only one that can annihilate a $y$) and one $\lambda_2$ vertex (to build up the remaining $x$ quanta via $x^2 z$).  The $z$ mode plays the role of the internal (virtual) line:

\begin{center}
  \begin{tikzpicture}
    \begin{feynman}
      \vertex (v1);
      \vertex [left=2cm of v1] (yin) {\(y\)};
      \vertex [above right=1.5cm of v1] (xout1) {\(x\)};
      \vertex [right=2cm of v1] (v2);
      \vertex [above right=1.5cm of v2] (xout2) {\(x\)};
      \vertex [below right=1.5cm of v2] (xout3) {\(x\)};
      \diagram*{
        (yin)   -- [scalar] (v1),
        (v1)    -- [scalar, edge label=\(x\)] (xout1),
        (v1)    -- [scalar, edge label'=\(z\) (virtual)] (v2),
        (v2)    -- [scalar] (xout2),
        (v2)    -- [scalar] (xout3),
      };
    \end{feynman}
  \end{tikzpicture}
\end{center}

\begin{itemize}
  \item Vertex $v_1$: corresponds to $\lambda_1 xyz$.  The $y$ is annihilated; one $x$ and one $z$ emerge.
  \item Internal line: the $z$ particle propagates from $v_1$ to $v_2$.  It is \emph{virtual} — its energy need not equal $\omega_z$.
  \item Vertex $v_2$: corresponds to $\lambda_2 x^2 z$.  The $z$ is annihilated; two $x$ particles emerge.
\end{itemize}

There is a second diagram corresponding to the other ordering (acting with $\lambda_2$ first, then $\lambda_1$), which gives the second term in \eqref{eq:lec04:amp-factored}.

\subsubsection*{From diagram to amplitude — preview of Feynman rules}

Each element of the Feynman diagram maps to a mathematical factor.  The full set of \textbf{Feynman rules} (to be derived in Lecture~05) will specify:
\begin{center}
\begin{tabular}{ll}
  \toprule
  Diagram element & Mathematical factor \\
  \midrule
  External line (particle $q$) & $\dfrac{1}{\sqrt{2\omega_q}}$ \\[8pt]
  Internal line (particle $q$) & $\dfrac{1}{E_i - E_q}$ (energy denominator) \\[8pt]
  Vertex with coupling $\lambda$ & $\lambda$ \\[4pt]
  Symmetry factor & $S$ (combinatorial, diagram-dependent) \\
  \bottomrule
\end{tabular}
\end{center}

Multiplying these factors together and summing all topologically distinct diagrams reproduces the perturbation-theory result — in this case \eqref{eq:lec04:amp-final} — without ever writing down an explicit sum over intermediate states.

%%──────────────────────────────────────────────────────────────────
\subsection*{To Remember}

\begin{itemize}
  \item \textbf{Fermi's Golden Rule}: $W \propto |\mathcal{A}|^2 \times \text{phase space}$.  The amplitude $|\mathcal{A}|$ carries all dynamical information; phase space carries kinematics.
  \item For a coupling $\alpha$ at a single vertex, the decay rate goes as $\Gamma \propto \alpha^2$ and the lifetime as $\tau \propto 1/\alpha^2$.
  \item \textbf{Selection rules} from $H_1 \propto x^p y^q$: only transitions with $|\Delta n_x| \leq p$ (same parity as $p$) and $|\Delta n_y| \leq q$ (same parity as $q$) are allowed.
  \item \textbf{Second-order perturbation theory}: $\mathcal{A}^{(2)} = \sum_n \langle f|H_1|n\rangle\langle n|H_1|i\rangle/(E_i-E_n)$.  Sum over \emph{all} intermediate states reachable from $|i\rangle$ and from which $|f\rangle$ can be reached.
  \item \textbf{Virtual (intermediate) states} need not conserve energy; only external states must.
  \item A \textbf{Feynman diagram} has external lines (initial/final), vertices (couplings), and internal lines (propagators/virtual particles).
  \item The \textbf{Feynman rules} (next lecture) translate each diagram element into a factor; the product of all factors gives the amplitude without explicit sums over intermediate states.
\end{itemize}

%%──────────────────────────────────────────────────────────────────
\subsection*{Further Reading}

No specific textbook was cited in this lecture.  The following are recommended:

\begin{itemize}
  \item \textbf{Griffiths, Chapter~6.1} — ``The Theory of Ionization and Excitation
    by Impact'': derives Fermi's Golden Rule in the particle-physics context and
    explains the role of the matrix element and the phase-space factor.  This is
    the closest Griffiths chapter to the perturbation-theory content of this
    lecture.

  \item \textbf{Peskin \& Schroeder, \S4.1--4.3} — Interaction picture, time-ordered
    products, and Wick's theorem.  These sections provide the formal QFT justification
    for the diagrammatic rules that were here derived for the harmonic oscillator and
    extended heuristically to field theory in Lecture~05.

  \item \textbf{Tong, \textit{QFT Notes}, Chapter~3, \S3.1--3.2}
    (\href{https://www.damtp.cam.ac.uk/user/tong/qft.html}{\texttt{damtp.cam.ac.uk/user/tong/qft.html}}):
    ``Interacting Fields'' — introduces the interaction picture and perturbation
    expansion in a QFT context, culminating in the same diagrammatic logic.

  \item \textbf{Sakurai \& Napolitano, \textit{Modern Quantum Mechanics}, Chapter~5}
    — Time-dependent perturbation theory and Fermi's Golden Rule in the
    non-relativistic QM setting.  A useful review of the underlying formalism
    before the field-theory generalisation.

  \item \textbf{Feynman, \textit{QED: The Strange Theory of Light and Matter}}
    (\href{https://press.princeton.edu/books/paperback/9780691164090/qed}{\texttt{Princeton UP}}):
    a non-technical but surprisingly deep account of how Feynman diagrams encode
    quantum amplitudes.  Recommended for conceptual grounding.
\end{itemize}
